%% ============================================================
%% BÖLÜM 13 — Güven Aralıkları
%% Olasılık Teorisi ve Matematiksel İstatistik
%% ============================================================
\chapter{Güven Aralıkları}
\label{chp:guven_araliklari}

\bolumalinti{%
  Statistics is the grammar of science.
}{Karl Pearson}

\begin{kazanimlar}
Bu bölümün sonunda okuyucu:
\begin{itemize}
  \item Güven aralığı tanımını, kapsama olasılığını ve doğru yorumunu açıklar.
  \item Pivot yöntemini kullanarak güven aralığı türetir.
  \item Normal, $t$, $\chi^2$ ve $F$ dağılımlarına dayalı güven aralıkları kurar.
  \item Büyük örneklem güven aralıklarını (Wald, oran, varyans stabilizasyonu) uygular.
  \item Aralık uzunluğu ile kapsama olasılığı arasındaki dengeyi analiz eder.
  \item Homojen olmayan varyanslar için Welch güven aralığını kullanır.
\end{itemize}
\end{kazanimlar}

\begin{motivasyon}
Nokta tahmini tek bir sayı verir; ama bu sayının ne kadar güvenilir olduğunu söylemez. Güven aralığı, belirsizliği nicelleştiren bir aralık tahminidir. Bir anket şirketi "\%43 ± \%2" derken aslında güven aralığı bildirmektedir.

\medskip
\noindent\textbf{Köprü:} Bölüm~\ref{chp:nokta_tahmin}'deki MLE'nin asimptotik normalliği, büyük örneklem güven aralıklarının temelidir. Bölüm~\ref{chp:dagilim_aileleri}'deki $t$, $\chi^2$, $F$ dağılımları, küçük örneklem güven aralıklarında kullanılacak.
\end{motivasyon}

%% ============================================================
\section{Güven Aralığı Kavramı}
\label{sec:ga_kavram}
%% ============================================================

\begin{tanim}[Güven aralığı]\label{def:guven_araligi}
$X_1,\ldots,X_n\sim F_\theta$ ve $L=L(\mathbf{X})$, $U=U(\mathbf{X})$ istatistikleri olsun ($L\leq U$). $(1-\alpha)$ \textbf{kapsama olasılıklı güven aralığı} \emph{(confidence interval, CI)}:
\[
  \Prob_\theta(L \leq \theta \leq U) \geq 1-\alpha \quad \text{her } \theta\in\Theta \text{ için.}
\]
$1-\alpha$ \textbf{güven düzeyi} \emph{(confidence level)}, $\alpha$ \textbf{anlamlılık düzeyi} \emph{(significance level)} adını alır.
\end{tanim}

\begin{tanimgerekce}
Sık yapılan yorum hatası: "\%95 olasılıkla $\theta$, $[L,U]$ içinde" — bu yanlış. $\theta$ sabit; $L$ ve $U$ rasgele. Doğru yorum: Yöntemi tekrar tekrar uygulandığında aralıkların \%95'i gerçek $\theta$'yı kapsar.
\end{tanimgerekce}

\begin{hataanalizi}
"Bu özel aralık $[1.2, 2.8]$'in \%95 olasılıkla $\mu$'yü kapsadığı" ifadesi hatalıdır. Gözlemlendikten sonra aralık sabittir; $\mu$ ya içindedir ya değildir. "Bu yöntemle üretilen aralıkların \%95'i $\mu$'yü kapsar" doğru ifadedir.
\end{hataanalizi}

\subsection{Pivot Yöntemi}

\begin{tanim}[Pivot]\label{def:pivot}
$Q = Q(X_1,\ldots,X_n;\theta)$ istatistiği, dağılımı $\theta$'dan bağımsız ise \textbf{pivot} \emph{(pivotal quantity)} adını alır.
\end{tanim}

\begin{teorem}[Pivot güven aralığı]\label{thm:pivot_ga}
$Q$ pivot ve $\Prob(a\leq Q\leq b)=1-\alpha$ ise, $a\leq Q\leq b$ eşitsizliğini $\theta$ için çözmek $(1-\alpha)$ güven aralığı verir.
\end{teorem}

%% ============================================================
\section{Normal Popülasyon: Bilinen Varyans}
\label{sec:normal_bilinen}
%% ============================================================

\begin{teorem}[$z$-güven aralığı]\label{thm:z_ga}
$X_1,\ldots,X_n\iid\Normal(\mu,\sigma^2)$, $\sigma^2$ bilinen. O zaman
\[
  Z = \frac{\overline{X}_n-\mu}{\sigma/\sqrt{n}} \sim \Normal(0,1)
\]
bir pivot. $(1-\alpha)$ güven aralığı:
\[
  \left(\overline{X}_n - z_{\alpha/2}\frac{\sigma}{\sqrt{n}},\; \overline{X}_n + z_{\alpha/2}\frac{\sigma}{\sqrt{n}}\right),
\]
$z_{\alpha/2} = \Phi^{-1}(1-\alpha/2)$ ile.
\end{teorem}

\begin{proof}
$\Prob(-z_{\alpha/2}\leq Z\leq z_{\alpha/2})=1-\alpha$. İçi çözersek $\overline{X}_n - z_{\alpha/2}\sigma/\sqrt{n}\leq\mu\leq\overline{X}_n+z_{\alpha/2}\sigma/\sqrt{n}$.
\end{proof}

\begin{ornek}[$z$-aralığı]\label{ex:z_aralik}
$n=25$, $\bar{x}=12.4$, $\sigma=2$ (bilinen). $\alpha=0.05$; $z_{0.025}=1.960$. GA: $12.4\pm1.960\cdot(2/5) = 12.4\pm0.784 = [11.616, 13.184]$.
\end{ornek}

%% ============================================================
\section{Normal Popülasyon: Bilinmeyen Varyans}
\label{sec:normal_bilinmeyen}
%% ============================================================

\begin{teorem}[$t$-güven aralığı]\label{thm:t_ga}
$X_1,\ldots,X_n\iid\Normal(\mu,\sigma^2)$, $\sigma^2$ bilinmiyor. O zaman
\[
  T = \frac{\overline{X}_n-\mu}{S/\sqrt{n}} \sim t_{n-1}
\]
bir pivot. $(1-\alpha)$ güven aralığı:
\[
  \left(\overline{X}_n - t_{n-1,\alpha/2}\frac{S}{\sqrt{n}},\; \overline{X}_n + t_{n-1,\alpha/2}\frac{S}{\sqrt{n}}\right).
\]
\end{teorem}

\begin{proof}
Bölüm~\ref{chp:dagilim_aileleri}'den: Normal popülasyonda $\overline{X}_n\perp S^2$, $\overline{X}_n\sim\Normal(\mu,\sigma^2/n)$, $(n-1)S^2/\sigma^2\sim\chi^2_{n-1}$. Dolayısıyla $T = Z/\sqrt{\chi^2_{n-1}/(n-1)}\sim t_{n-1}$.
\end{proof}

\begin{ornek}[$t$-aralığı]\label{ex:t_aralik}
$n=16$, $\bar{x}=5.2$, $s=1.8$, $\sigma^2$ bilinmiyor, $\alpha=0.05$. $t_{15,0.025}=2.131$. GA: $5.2\pm2.131\cdot(1.8/4) = 5.2\pm0.959 = [4.241, 6.159]$.
\end{ornek}

\begin{kritiksorgu}
$n$ büyüdükçe $t$-aralığı $z$-aralığına yaklaşır ($t_{n-1}\to\Normal(0,1)$). Pratikte $n\geq30$ için fark ihmal edilebilir; ama normallik varsayımı hâlâ önemlidir.
\end{kritiksorgu}

%% ============================================================
\section{Varyans için Güven Aralığı}
\label{sec:varyans_ga}
%% ============================================================

\begin{teorem}[$\chi^2$-güven aralığı]\label{thm:chi2_ga}
$X_1,\ldots,X_n\iid\Normal(\mu,\sigma^2)$. O zaman
\[
  Q = \frac{(n-1)S^2}{\sigma^2} \sim \chi^2_{n-1}
\]
bir pivot. $(1-\alpha)$ güven aralığı:
\[
  \left(\frac{(n-1)S^2}{\chi^2_{n-1,\alpha/2}},\; \frac{(n-1)S^2}{\chi^2_{n-1,1-\alpha/2}}\right).
\]
\end{teorem}

\begin{proof}
$\Prob(\chi^2_{n-1,1-\alpha/2}\leq Q\leq\chi^2_{n-1,\alpha/2})=1-\alpha$; $\sigma^2$ için çöz.
\end{proof}

\begin{hataanalizi}
$\chi^2$ aralığı normallik varsayımına son derece duyarlıdır. Dağılım normalden sapıyorsa bootstrap yöntemi tercih edilir.
\end{hataanalizi}

%% ============================================================
\section{İki Örneklem Karşılaştırması}
\label{sec:iki_orneklem}
%% ============================================================

\subsection{Eşit Varyanslı İki Ortalama}

\begin{teorem}[İki örneklem $t$-aralığı]\label{thm:iki_t}
$X_1,\ldots,X_m\iid\Normal(\mu_1,\sigma^2)$ ve $Y_1,\ldots,Y_n\iid\Normal(\mu_2,\sigma^2)$ bağımsız (ortak $\sigma^2$). Havuzlanmış varyans tahmin edicisi:
\[
  S_p^2 = \frac{(m-1)S_X^2+(n-1)S_Y^2}{m+n-2}.
\]
$\mu_1-\mu_2$ için $(1-\alpha)$ güven aralığı:
\[
  (\bar{X}-\bar{Y}) \pm t_{m+n-2,\alpha/2}\,S_p\!\sqrt{\frac{1}{m}+\frac{1}{n}}.
\]
\end{teorem}

\begin{proof}
$\bar{X}-\bar{Y}\sim\Normal(\mu_1-\mu_2,\sigma^2(1/m+1/n))$ ve $(m+n-2)S_p^2/\sigma^2\sim\chi^2_{m+n-2}$ bağımsız olduğundan pivot $t_{m+n-2}$'dir.
\end{proof}

\subsection{Welch Güven Aralığı (Eşit Olmayan Varyanslarda)}

\begin{teorem}[Welch $t$-aralığı]\label{thm:welch}
$\sigma_1^2\neq\sigma_2^2$ (bilinmiyor). Satterthwaite serbestlik derecesi:
\[
  \nu = \frac{(S_X^2/m + S_Y^2/n)^2}{\dfrac{(S_X^2/m)^2}{m-1}+\dfrac{(S_Y^2/n)^2}{n-1}}.
\]
$\mu_1-\mu_2$ için yaklaşık güven aralığı:
\[
  (\bar{X}-\bar{Y}) \pm t_{\nu,\alpha/2}\!\sqrt{\frac{S_X^2}{m}+\frac{S_Y^2}{n}}.
\]
\end{teorem}

\begin{ispatiskele}
Welch-Satterthwaite yaklaşımı, $S_X^2/m+S_Y^2/n$'yi $a\chi^2_\nu$ ile eşleştirir; $a$ ve $\nu$ moment eşleştirmesiyle belirlenir. Yaklaşık (asimptotik değil); normallik gerekir.
\end{ispatiskele}

%% ============================================================
\section{Büyük Örneklem Güven Aralıkları}
\label{sec:buyuk_orneklem}
%% ============================================================

\begin{teorem}[Wald güven aralığı]\label{thm:wald_ga}
Regularity koşulları altında $\sqrt{n}(\MLE\theta-\theta)\dto\Normal(0,1/\mathcal{I}(\theta))$. Büyük örneklem $(1-\alpha)$ güven aralığı:
\[
  \MLE\theta \pm \frac{z_{\alpha/2}}{\sqrt{n\,\hat{\mathcal{I}}(\MLE\theta)}},
\]
$\hat{\mathcal{I}}$ gözlemlenen Fisher bilgisi.
\end{teorem}

\begin{ornek}[Bernoulli oranı için GA]\label{ex:oran_ga}
$X_1,\ldots,X_n\iid\mathrm{Ber}(p)$. $\MLE{p}=\bar{X}$, $\mathcal{I}(p)=1/(p(1-p))$. Wald aralığı:
\[
  \bar{X} \pm z_{\alpha/2}\sqrt{\frac{\bar{X}(1-\bar{X})}{n}}.
\]
\end{ornek}

\begin{hataanalizi}
Wald oranı güven aralığı $p\approx0$ veya $p\approx1$'de ya da küçük $n$'de kötü performans gösterir. Wilson aralığı (skorlu GA) daha sağlamdır:
\[
  \frac{\bar{X}+z_{\alpha/2}^2/(2n)}{1+z_{\alpha/2}^2/n} \pm \frac{z_{\alpha/2}\sqrt{\bar{X}(1-\bar{X})/n + z_{\alpha/2}^2/(4n^2)}}{1+z_{\alpha/2}^2/n}.
\]
\end{hataanalizi}

\subsection{Varyans Stabilizasyon Dönüşümü}

\begin{teorem}[Delta yöntemi ile GA]\label{thm:delta_ga}
$g$ türevlenebilir ve $g'(\theta)\neq0$. Delta yönteminden $\sqrt{n}(g(\MLE\theta)-g(\theta))\dto\Normal(0,[g'(\theta)]^2/\mathcal{I}(\theta))$. $g$ seçimi ile varyans sabitlenirse, GA sabit genişlikte olur.
\end{teorem}

\begin{ornek}[Poisson varyans stabilizasyonu]
$\bar{X}\sim\Pois(\lambda)$: $\Var(\bar{X})\approx\lambda/n$. $g(\lambda)=\sqrt{\lambda}$ ile delta yöntemi: $\sqrt{n}(\sqrt{\bar{X}}-\sqrt{\lambda})\dto\Normal(0,1/4)$. GA:
\[
  \sqrt{\bar{X}} \pm \frac{z_{\alpha/2}}{2\sqrt{n}} \quad\Rightarrow\quad
  \lambda: \left(\sqrt{\bar{X}} - \frac{z_{\alpha/2}}{2\sqrt{n}}\right)^{\!2} \leq \lambda \leq \left(\sqrt{\bar{X}} + \frac{z_{\alpha/2}}{2\sqrt{n}}\right)^{\!2}.
\]
\end{ornek}

%% ============================================================
\section{Güven Aralığı Uzunluğu ve Örneklem Büyüklüğü}
\label{sec:ga_uzunluk}
%% ============================================================

\begin{teorem}[Örneklem büyüklüğü belirleme]\label{thm:orneklem_buyuklugu}
$z$-aralığı ile $\mu$'yü $\pm d$ hassasiyetle tahmin etmek için gereken minimum örneklem büyüklüğü:
\[
  n \geq \left\lceil \frac{z_{\alpha/2}^2\,\sigma^2}{d^2} \right\rceil.
\]
$\sigma$ bilinmiyorsa pilot çalışma veya muhafazakâr $\sigma$ tahmini kullanılır.
\end{teorem}

\begin{proof}
GA yarı-genişliği $z_{\alpha/2}\sigma/\sqrt{n}\leq d$ eşitsizliğini $n$ için çöz.
\end{proof}

\begin{ornek}[Seçim anketi]
Bir oy oranını $\pm0.03$ hassasiyetle \%95 güvenle tahmin et. Muhafazakâr $p=0.5$ ile $\sigma^2=p(1-p)=0.25$. $n\geq\lceil 1.96^2\cdot0.25/0.03^2\rceil = \lceil 1067.1\rceil = 1068$.
\end{ornek}

\begin{kendinleifade}
Güven düzeyini $\alpha=0.05$'ten $\alpha=0.01$'e çekince aralık nasıl değişir? Örneklem büyüklüğünü 4 katına çıkarmak aralık genişliğini kaçta bire indirir?
\kendinleifadecevap{Güven düzeyi $\alpha=0.05\to0.01$: $z_{0.025}=1.96\to z_{0.005}=2.576$,
  aralık genişler (daha güvenilir = daha geniş).
  Örneklem $n\to4n$: Genişlik $\propto\sigma/\sqrt{n}\to\sigma/\sqrt{4n}=\tfrac{1}{2}\cdot\sigma/\sqrt{n}$,
  yarıya iner. Sezgi: güveni artırmak genişletir, veriyi artırmak daraltır;
  ikisi birbirini dengeler.}
\end{kendinleifade}

%% ============================================================
\section{En Kısa Güven Aralığı}
\label{sec:en_kisa_ga}
%% ============================================================

Birden fazla güven aralığı aynı kapsama olasılığına sahip olabilir. En bilgilisi en kısasıdır.

\begin{teorem}[Simetrik unimodal dağılımda en kısa GA]\label{thm:en_kisa}
Pivot $Q$'nun yoğunluğu unimodal ve simetrik ($f_Q$ genişleme merkezi $c$ olan) ise, sabit uzunluklu $(1-\alpha)$ aralık içinde kapsama olasılığını maksimize eden — eşdeğer olarak, sabit kapsama olasılığında uzunluğu minimize eden — aralık simetriktir: $\Prob(c-d\leq Q\leq c+d)=1-\alpha$.
\end{teorem}

\begin{proof}[Proof (taslak)]
Neyman-Pearson argümanı: $f_Q(q)>k$ kümesi üzerinde topla; $k$ normalleştirme sabitine göre ayarla. Unimodal simetrik $f_Q$ için bu küme simetrik bir aralık.
\end{proof}

\begin{sonuc}[Normal güven aralığı optimal]\label{cor:normal_optimal}
$\Normal(0,1)$ pivot için simetrik $z$-aralığı, verilen $1-\alpha$ değerinde en kısa olasılıklı güven aralığıdır.
\end{sonuc}

\begin{ornek}[$\chi^2$ en kısa aralığı]\label{ex:chi2_en_kisa}
$\chi^2_{19}$ asimetrik olduğundan simetrik eşit kuyruk aralığı en kısa değildir. Optimal kesim noktaları $c_L, c_U$ için:
\[
  f_{\chi^2_{19}}(c_L) = f_{\chi^2_{19}}(c_U), \quad \Prob(c_L\leq Q\leq c_U)=0.95.
\]
Sayısal çözüm gerektirir; pratikte eşit kuyruk aralığı standart olarak kullanılır.
\end{ornek}

%% ============================================================
\section{Profil Olabilirlik Güven Aralığı}
\label{sec:profil_ga}
%% ============================================================

\begin{tanim}[Profil olabilirlik]\label{def:profil_likh}
$\theta=(\psi,\lambda)$, ilgi parametresi $\psi$, tortu $\lambda$. \textbf{Profil olabilirlik}:
\[
  L_P(\psi) = \max_{\lambda} L(\psi,\lambda) = L(\psi,\hat\lambda(\psi)).
\]
\textbf{Profil olabilirlik oran GA} (Wilks teoreminden):
\[
  R_P(\psi) = 2[\ell(\hat\theta) - \ell_P(\psi)] \leq \chi^2_{1,\alpha}.
\]
\end{tanim}

\begin{teorem}[Wilks — profil olabilirlik GA]\label{thm:wilks_profil}
Düzenli model, $n\to\infty$. O zaman $\{\psi: R_P(\psi)\leq\chi^2_{1,\alpha}\}$ yaklaşık $(1-\alpha)$ güven aralığı oluşturur.
\end{teorem}

\begin{ornek}[Lognormal parametre aralığı]\label{ex:profil_lognormal}
$X_i\iid\mathrm{LogNormal}(\mu,\sigma^2)$. $\psi=e^\mu$ ilgi parametresi. Profil GA:
1. $\hat\mu,\hat\sigma^2$ = MLE. 2. Farklı $\mu$ değerleri için $\hat\sigma^2(\mu)=\frac1n\sum(\log X_i-\mu)^2$. 3. $R_P(\mu)\leq3.84$ kümesi GA.

Profil GA, Wald GA'ya göre asimetrik olabilir — gerçek GA şekline daha iyi uyar.
\end{ornek}

%% ============================================================
\section{Eşzamanlı Güven Aralıkları}
\label{sec:esz_ga}
%% ============================================================

\begin{tanim}[Eşzamanlı GA]\label{def:esz_ga}
$k$ parametre $(\theta_1,\ldots,\theta_k)$ için aralıklar $(L_j,U_j)$, $j=1,\ldots,k$. \textbf{Aile genel hata oranı} (FWER):
\[
  \Prob(\exists j: \theta_j\notin[L_j,U_j]) \leq \alpha.
\]
Bireysel $\alpha_j=\alpha/k$ (Bonferroni) bunu sağlar ama muhafazakârdır.
\end{tanim}

\begin{teorem}[Bonferroni düzeltmesi]\label{thm:bonferroni_ga}
Her aralık $\alpha/k$ düzeyinde kurulursa FWER $\leq\alpha$.
\end{teorem}

\begin{proof}
Birleşim sınırı: $\Prob(\bigcup_j A_j)\leq\sum_j\Prob(A_j)\leq k\cdot(\alpha/k)=\alpha$.
\end{proof}

\begin{teorem}[Scheffé eşzamanlı GA — regresyon]\label{thm:scheffe_ga}
Çok değişkenli regresyonda tüm doğrusal kombinasyonlar $c^\top\beta$ için eşzamanlı \%$95$ GA:
\[
  c^\top\hat\beta \pm \sqrt{p\,F_{p,n-p,\alpha}}\cdot\hat\sigma\sqrt{c^\top(X^\top X)^{-1}c}.
\]
Bu, sonsuz sayıda kombinasyonu eşzamanlı kapsar; Bonferroni'den geniş ama kapsamlıdır.
\end{teorem}

%% ============================================================
\section{Güven Bölgeleri (Çok Parametreli)}
\label{sec:guven_bolgesi}
%% ============================================================

\begin{tanim}[Güven bölgesi]\label{def:guven_bolgesi}
$p$ boyutlu parametre $\boldsymbol\theta$. Rasgele küme $C(\mathbf{X})\subseteq\mathbb{R}^p$ için $(1-\alpha)$ \textbf{güven bölgesi} \emph{(confidence region)}:
\[
  \Prob_{\boldsymbol\theta}(\boldsymbol\theta\in C(\mathbf{X})) \geq 1-\alpha \quad \text{her } \boldsymbol\theta\in\Theta \text{ için.}
\]
\end{tanim}

\begin{teorem}[Eliptik güven bölgesi]\label{thm:eliptik_bolge}
$X_1,\ldots,X_n\iid\Normal(\boldsymbol\mu,\Sigma)$, $\Sigma$ bilinen. $(\bar{\mathbf{X}}-\boldsymbol\mu)^\top (n\Sigma^{-1})(\bar{\mathbf{X}}-\boldsymbol\mu)\sim\chi^2_p$ pivottur. $(1-\alpha)$ güven bölgesi:
\[
  C(\mathbf{X}) = \{\boldsymbol\mu: (\bar{\mathbf{X}}-\boldsymbol\mu)^\top (n\Sigma^{-1})(\bar{\mathbf{X}}-\boldsymbol\mu) \leq \chi^2_{p,\alpha}\}.
\]
Bu bir $p$-boyutlu elipsoittir; $\bar{\mathbf{X}}$ merkezi, $\Sigma/(n\chi^2_{p,\alpha})$ yarı-eksen uzunlukları.
\end{teorem}

\begin{proof}
$\bar{\mathbf{X}}\sim\Normal(\boldsymbol\mu,\Sigma/n)$; standartlaştırarak $\mathbf{Z}=\Sigma^{-1/2}(\bar{\mathbf{X}}-\boldsymbol\mu)\sqrt{n}\sim\Normal(\mathbf{0},I_p)$. $\|\mathbf{Z}\|^2=(\bar{\mathbf{X}}-\boldsymbol\mu)^\top(n\Sigma^{-1})(\bar{\mathbf{X}}-\boldsymbol\mu)\sim\chi^2_p$ (standart Normal karelerinin toplamı).
\end{proof}

\begin{teorem}[Asimptotik güven bölgesi — genel]\label{thm:asimptotik_bolge}
MLE $\hat{\boldsymbol\theta}_n$ için:
\[
  n(\hat{\boldsymbol\theta}_n-\boldsymbol\theta)^\top I(\boldsymbol\theta) (\hat{\boldsymbol\theta}_n-\boldsymbol\theta) \dto \chi^2_p.
\]
Büyük örneklem güven bölgesi: $\{\boldsymbol\theta: n(\hat{\boldsymbol\theta}_n-\boldsymbol\theta)^\top \hat{I}(\hat{\boldsymbol\theta}_n)(\hat{\boldsymbol\theta}_n-\boldsymbol\theta)\leq\chi^2_{p,\alpha}\}$.
\end{teorem}

\begin{ornek}[İki parametreli Normal güven elipsi]
$X_1,\ldots,X_n\iid\Normal(\mu,\sigma^2)$, ikisi bilinmiyor; ilgi parametresi $(\mu,\sigma^2)$. Fisher bilgisi diyagonal: $I_{11}=1/\sigma^2$, $I_{22}=1/(2\sigma^4)$. Elipsin eksenler $(\mu-\hat\mu)^2/(n\hat\sigma^2)$ ve $(\sigma^2-\hat\sigma^2)^2/(2\hat\sigma^4/n)$ yönünde. İki boyutlu elipsin alanı tek boyutlu aralıkların yüzey alanıyla değil yüzey dışıyla kıyaslanır (Bonferroni daha muhafazakâr).
\end{ornek}

\begin{kritiksorgu}
Bireysel $p$ adet \%95 GA'yı oluşturmak yerine neden doğrudan güven bölgesi kurarız? Bireysel aralıkların kesişimi güven bölgesi değildir: iki \%95 aralığın kesişimi en az \%90 kapsama garantisi verir (Bonferroni). Eliptik güven bölgesi tam $1-\alpha$ garantisiyle daha verimli bilgi taşır.
\end{kritiksorgu}

%% ============================================================
\section{Tolerans Aralıkları}
\label{sec:tolerans_aralik}
%% ============================================================

\begin{tanim}[Tolerans aralığı]\label{def:tolerans_aralik}
$(1-\alpha)$ güvenle popülasyonun en az $\gamma$ oranını kapsayan aralık \textbf{tolerans aralığı} \emph{(tolerance interval)}: $L,U$ rasgeleler; gözlemlenmemiş $X_{\text{yeni}}$ için
\[
  \Prob\!\bigl[\Prob(L\leq X_{\text{yeni}}\leq U\mid\mathbf{X})\geq\gamma\bigr]\geq1-\alpha.
\]
Tolerans aralığı, güven aralığından farklıdır: bireysel gözlemleri kapsar, parametreyi değil.
\end{tanim}

\begin{teorem}[Normal tolerans aralığı]\label{thm:normal_tolerans}
$X_1,\ldots,X_n\iid\Normal(\mu,\sigma^2)$; $\mu,\sigma^2$ bilinmiyor. $(1-\alpha)$ güven ve $\gamma$ kapsama oranı için tolerans aralığı:
\[
  \bar{X} \pm k\,S,
\]
$k = k(n,\gamma,\alpha)$ tolerans çarpanı tablodan veya şu yaklaşımla: $k\approx\frac{z_{(1-\gamma)/2}}{\sqrt{1-z_\alpha^2/(2(n-1))}}\sqrt{1+\frac{1}{n}}$ (Wald-Wolfowitz yaklaşımı).
\end{teorem}

\begin{proof}[Proof — taslak]
$X_{\text{yeni}}|\mu,\sigma\sim\Normal(\mu,\sigma^2)$; $\bar X|\mu,\sigma\sim\Normal(\mu,\sigma^2/n)$. $\Prob(|\bar X-X_{\text{yeni}}|>kS)$ ifadesi $\chi^2$ ve $t$ dağılımlarını içerir; integral kapalı formda değil. $k$ tablo değerleri Monte Carlo veya yaklaşımla hesaplanır.
\end{proof}

\begin{ornek}[Tolerans aralığı — vida uzunluğu]
Makine $n=25$ vida üretti; $\bar{x}=50.2$ mm, $s=0.3$ mm. \%95 güvenle vidaların \%99'unu kapsayan tolerans aralığını bul. $k(25,0.99,0.05)\approx3.46$ (tablodan). Aralık: $50.2\pm3.46\cdot0.3=50.2\pm1.04=[49.16,\;51.24]$ mm.
\end{ornek}

\begin{hataanalizi}
Üç farklı aralık türünü karıştırmayın:
\begin{itemize}
  \item \textbf{Güven aralığı}: $\mu$ gibi bir parametre için — aralık ortalamayı kapsar.
  \item \textbf{Tahmin aralığı}: tek bir yeni gözlem için — $X_{\text{yeni}}\in[L,U]$.
  \item \textbf{Tolerans aralığı}: popülasyonun $\gamma$ oranı için — hem gelecek hem geçmiş gözlemler.
\end{itemize}
Tolerans aralığı en geniş, güven aralığı en dar. Kalite güvencesinde tolerans aralığı standart.
\end{hataanalizi}

\begin{teorem}[Tahmin aralığı — tek yeni gözlem]\label{thm:tahmin_aralik}
$X_{\text{yeni}}$ bağımsız, $\Normal(\mu,\sigma^2)$'dan. $\mu,\sigma^2$ bilinmiyor. $(1-\alpha)$ tahmin aralığı \emph{(prediction interval)}:
\[
  \bar X \pm t_{n-1,\alpha/2}\,S\sqrt{1+\frac{1}{n}}.
\]
Güven aralığından $\sqrt{1+1/n}$ faktörü kadar geniş (yeni gözlem belirsizliği ekstra).
\end{teorem}

\begin{disiplinlerarasibaglan}
\textbf{Kalite güvencesi:} Tolerans aralıkları üretim spesifikasyonu belirleme standardı (ISO 16269-6). \textbf{Tıbbi teşhis:} Referans aralıkları (sağlıklı popülasyonun \%95'ini kapsayan değerler) birer tolerans aralığıdır. \textbf{Çevre mühendisliği:} Kirletici konsantrasyonu için üst tolerans sınırı yasal düzenleme referansı.
\end{disiplinlerarasibaglan}

%% ============================================================
\section{Güven Aralıklarında Asimptotik Doğruluk: Edgeworth Düzeltmesi}
\label{sec:edgeworth_ga}
%% ============================================================

Standart Wald güven aralığı $\bar X_n \pm z_{\alpha/2}\sigma/\sqrt{n}$ adı altında kapsama olasılığı $\Phi(z_{\alpha/2})-\Phi(-z_{\alpha/2})$'ye yaklaşır; ancak bu yaklaşımın hata mertebesi $O(n^{-1/2})$'dir. Bartlett ve diğerleri, Edgeworth açılımı kullanarak daha kesin kapsama olasılığı hesaplamıştır.

\begin{teorem}[Edgeworth düzeltmesi]\label{thm:edgeworth_ga}
$X_i$ i.i.d., sonlu 4. momenti var, $\gamma_1=\mu_3/\sigma^3$ (çarpıklık). Wald güven aralığının gerçek kapsama olasılığı:
\[
  P\!\left(\theta\in\bar X_n\pm\frac{z_{\alpha/2}\sigma}{\sqrt{n}}\right)
  = (1-\alpha) + \frac{\gamma_1}{6\sqrt{n}}[z_{\alpha/2}^2-1]\cdot2\phi(z_{\alpha/2}) + O(n^{-1}).
\]
Çarpıklık terimi $\gamma_1\neq0$ ise kapsama olasılığını $(1-\alpha)$'dan saptırır.
\end{teorem}

\begin{proof}[Taslak]
$W_n=\sqrt{n}(\bar X_n-\mu)/\sigma$ için Edgeworth açılımı:
\[
  F_{W_n}(w) = \Phi(w) - \frac{\gamma_1}{6\sqrt{n}}(w^2-1)\phi(w) + O(n^{-1}).
\]
Kapsama olasılığını $F_{W_n}(z_{\alpha/2})-F_{W_n}(-z_{\alpha/2})$ olarak yazıp Edgeworth açılımını yerine koy; $\phi$ simetrik, $(w^2-1)$ anti-simetrik olmadığından fark birinci terimde sıfırlanmaz; teoremde verilen düzeltme elde edilir.
\end{proof}

\begin{ornek}[Üstel dağılımda kapsama hatası]
$X_i\sim\mathrm{Exp}(\lambda)$: $\gamma_1=2$. $n=20$, $1-\alpha=0.95$ ($z_{0.025}=1.96$).
Edgeworth düzeltmesi:
\[
  \Delta = \frac{2}{6\sqrt{20}}(1.96^2-1)\cdot2\phi(1.96) = \frac{2}{6\cdot4.47}(2.842)(0.584)\approx0.0124.
\]
Gerçek kapsama $\approx0.95+0.012=0.962$ (aşırı kapsama) — sağa çarpık dağılımda normal yaklaşım tutucu davranır.
\end{ornek}

\begin{teorem}[Bootstrap-t güven aralığının kapsama doğruluğu]\label{thm:bootstrap_kapsama}
Double bootstrap güven aralığı (Beran 1987), kapsama hatasını $O(n^{-1})$'den $O(n^{-3/2})$'ye indirger. Percentile-t bootstrap $O(n^{-1})$ doğruluk sağlar; basit yüzdelik bootstrap ise yalnızca $O(n^{-1/2})$.
\end{teorem}

\noindent\textbf{Sonuç:} Asimetrik dağılımlarda veya küçük $n$'de, double-bootstrap veya profil olabilirlik aralıkları, Wald aralığından belirgin şekilde üstündür.

%% ============================================================
\section{Adaptif Güven Aralıkları: Stein Küçültmesi ve SURE}
\label{sec:adaptif_ga}
%% ============================================================

$p\geq3$ parametreli $\boldsymbol\mu\in\mathbb{R}^p$ tahmininde, James-Stein tahminedici MLE'nin OKH'sini düşürür; bu fenomen güven bölgeleri için de paralel sonuçlar doğurur.

\begin{tanim}[SURE — Stein'ın hesapsız risk tahmincisi]
$\mathbf{X}\sim\Normal_p(\boldsymbol\mu,\sigma^2I)$, yansız risk tahmincisi:
\[
  \widehat{R}(\delta,\boldsymbol\mu) = p\sigma^2 + \sum_{i=1}^p\left[\left(\frac{\partial\delta_i}{\partial X_i}\right)^2\sigma^4 - 2\sigma^2\frac{\partial\delta_i}{\partial X_i}\right].
\]
Stein lemması ile beklenti $\boldsymbol\mu$'ya bağlı değildir; böylece risk veri gözlemlenerek tahmin edilebilir.
\end{tanim}

\begin{teorem}[James-Stein güven küresi]
$\mathbf{X}\sim\Normal_p(\boldsymbol\mu,I)$ ($p\geq3$), $\delta^{JS}=\left(1-\frac{p-2}{\|\mathbf{X}\|^2}\right)\mathbf{X}$ James-Stein tahminedici. Standart güven küresi $\{\boldsymbol\mu:\|\boldsymbol\mu-\mathbf{X}\|\leq c\}$'in kapsama olasılığı $(1-\alpha)$; James-Stein merkezli küre \emph{daha küçük} ama \emph{aynı ya da daha yüksek} kapsama sağlar (Stein 1981).
\end{teorem}

Bu, yüksek boyutlu istatistikte güven kümelerinin, doğal kayıp (Öteki yöntem: Bayes HPD bölgesi) ile tanımlandığında, sezgilerin aksine davranabileceğini gösterir.

%% ============================================================

\begin{mikroozet}
\begin{itemize}
  \item Güven aralığı: $\theta$ sabit, aralık rasgele; kapsama yorumu yönteme aittir, tek bir aralığa değil.
  \item Pivot yöntemi: dağılımı $\theta$'dan bağımsız büyüklük bul, eşitsizliği çöz.
  \item $z$/$t$/$\chi^2$ aralıkları normallik gerektirir; Welch eşitsiz varyanslarda güçlüdür.
  \item Wald büyük örneklem için evrensel; küçük oran ve uç değerlerde Wilson daha iyidir.
  \item En kısa GA: unimodal simetrik pivotta simetrik kesimler optimal.
  \item Profil olabilirlik GA: asimetrik ve parametrizasyona değişmez.
\end{itemize}
\end{mikroozet}

\begin{ozyansima}
Güven aralığının uzunluğu ne zaman belirsizliğimizi tam yansıtır? Güven aralığı ile Bayes güvenilir aralığı arasındaki felsefi fark nedir?
\end{ozyansima}

\begin{kopru}
\textbf{Nereden geldik:} Bölüm~\ref{chp:nokta_tahmin}'deki Fisher bilgisi ve asimptotik normallik Wald aralığını, Bölüm~\ref{chp:dagilim_aileleri}'deki $t/\chi^2/F$ dağılımları küçük örneklem aralıklarını mümkün kıldı.

\textbf{Nereye gidiyoruz:} Bölüm~\ref{chp:hipotez_testi}'de "bu $\theta_0$ değeri güven aralığında mı?" sorusu hipotez testine dönüşüyor. Güven aralığı ve hipotez testi dualdir.
\end{kopru}

%% ============================================================
%% ALISTIRMALAR
%% ============================================================

\cozumlualistirmalar

\begin{alistirma}[Al.13.1]{A}{$z$-aralığı}
$n=100$ gözlem; $\bar{x}=4.82$, $\sigma=1.2$ (bilinen). $\mu$ için \%90 ve \%99 güven aralıklarını hesapla ve uzunluklarını karşılaştır.
\end{alistirma}
\alcozum{%
$z_{0.05}=1.645$, $z_{0.005}=2.576$. \%90: $4.82\pm1.645\cdot0.12=[4.623,5.017]$ (uzunluk 0.394). \%99: $4.82\pm2.576\cdot0.12=[4.511,5.129]$ (uzunluk 0.617). Güven artınca uzunluk artar.%
}

\begin{alistirma}[Al.13.2]{B}{$t$-aralığı ve normallik}
$n=9$, $\bar{x}=23.1$, $s=4.5$, normal popülasyon. (a) \%95 $t$-aralığı. (b) $n=9$ yerine $n=100$ ile yaklaşık $z$-aralığını karşılaştır.
\end{alistirma}
\alcozum{%
(a) $t_{8,0.025}=2.306$; GA: $23.1\pm2.306\cdot(4.5/3)=23.1\pm3.459=[19.641,26.559]$. (b) $n=100$: $z_{0.025}=1.960$; GA: $23.1\pm1.960\cdot(4.5/10)=23.1\pm0.882$. Büyük örneklemde $t\approx z$ ve aralık çok daralır.%
}

\begin{alistirma}[Al.13.2b]{B}{Bootstrap Güven Aralığı}
$X_1,\ldots,X_{20}$ gözlemi, $\bar{x}=12.4$, $s=3.1$. (a) Nonparametrik bootstrap ile medyan için \%95 GA'nın mantığını açıklayın. (b) $B=1000$ bootstrap örnekleminde medyan standart hatası $\hat\sigma^*=0.82$ bulundu. Normal yaklaşım ile GA'yı kurun. (c) Bootstrap'ın klasik $t$-aralığına göre avantajı nedir?
\end{alistirma}
\alcozum{%
(a) Her bootstrap örnekleminden medyan $\tilde\theta^{*(b)}$ hesaplanır ($b=1,\ldots,B$). Emprik dağılım $F^*$ kullanılarak güven sınırları: $(\tilde\theta^*_{(\alpha/2)},\tilde\theta^*_{(1-\alpha/2)})$ yüzdelik yöntemi ile.
(b) Normal yaklaşım: $\tilde\theta\pm z_{\alpha/2}\hat\sigma^*=\tilde\theta\pm1.96\cdot0.82$. Medyan $\approx\bar{x}=12.4$ (simetrik ise): $12.4\pm1.607=[10.79,14.01]$.
(c) Bootstrap normallik varsaymaz; medyan, oran ve diğer düzgün olmayan istatistikler için geçerlidir; küçük örneklemde $t$-dağılımı varsayımını aşar.}

\begin{alistirma}[Al.13.2c]{B}{Bayes Güvenilir Aralığı — Hesap}
$X_1,\ldots,X_n\iid\Binom(1,p)$, $n=30$, $\sum x_i=18$. Önsel $p\sim\Beta(2,2)$. (a) Sonsal dağılımı bulun. (b) \%95 eşit-kuyruklu güvenilir aralığı Beta kantil tablosuyla kurun. (c) Bu aralığı klasik Wald aralığıyla karşılaştırın.
\end{alistirma}
\alcozum{%
(a) Sonsal: $p|X\sim\Beta(2+18,2+12)=\Beta(20,14)$.
(b) \%95 GA: $(\Beta_{0.025}(20,14), \Beta_{0.975}(20,14))$. $\Beta(20,14)$ modu $19/32=0.594$. Yaklaşık hesap: $\mu_1=20/34=0.588$, $\sigma_1^2=20\cdot14/(34^2\cdot35)=0.00696$. Normal yaklaşım: $0.588\pm1.96\cdot0.0834\approx[0.425,0.751]$.
(c) Wald: $\hat p=18/30=0.6$, $\hat\sigma=\sqrt{0.6\cdot0.4/30}=0.0894$; GA: $0.6\pm1.96\cdot0.0894=[0.425,0.775]$. Bayesçi önceden $\Beta(2,2)$ ortalamaya (0.5) doğru çeker.}

\alistirmalar

\begin{alistirma}[Al.13.3]{A}{Varyans için GA}
$n=20$, $s^2=16.4$, normal popülasyon. $\sigma^2$ için \%95 güven aralığını kur.
\end{alistirma}
\alcozum{%
$\chi^2_{19,0.025}=32.852$, $\chi^2_{19,0.975}=8.907$. GA: $[19\cdot16.4/32.852,\;19\cdot16.4/8.907]=[9.49,\;34.99]$.%
}

\begin{alistirma}[Al.13.4]{B}{Wilson oranı aralığı}
$n=20$ denemede $k=2$ başarı. $p$ için (a) Wald aralığı, (b) Wilson aralığı kur (\%95). Sonuçları yorumla.
\end{alistirma}
\alcozum{%
$\bar{p}=0.1$. (a) Wald: $0.1\pm1.96\sqrt{0.1\cdot0.9/20}=0.1\pm0.131=[-0.031,0.231]$ — olumsuz değer içeriyor, geçersiz. (b) Wilson: pay=$0.1+1.96^2/40=0.196$; payda=$1+3.84/20=1.192$; merkez$=0.164$; $\pm\approx0.155$; $[0.009,\;0.319]$ — daha sağlam.%
}

\begin{alistirma}[Al.13.5]{B}{Welch güven aralığı}
İki bağımsız grup: Grup A ($n_1=15$, $\bar{x}_1=42$, $s_1=8$), Grup B ($n_2=20$, $\bar{x}_2=37$, $s_2=12$). Varyansların eşit olmadığını varsayarak $\mu_1-\mu_2$ için \%95 Welch güven aralığını kur.
\end{alistirma}
\alcozum{%
SE $=\sqrt{64/15+144/20}=\sqrt{4.267+7.2}=\sqrt{11.467}=3.387$. $\nu=11.467^2/[(4.267^2/14)+(7.2^2/19)]=131.5/(1.298+2.737)=32.6\approx33$. $t_{33,0.025}\approx2.035$. GA: $5\pm2.035\cdot3.387=5\pm6.89=[-1.89,\;11.89]$.%
}

\begin{alistirma}[Al.13.6]{C}{Örneklem büyüklüğü optimizasyonu}
Bir üretim hattında $\sigma=5$ mm. Ortalamayı $\pm0.5$ mm hassasiyetle \%95 ile tahmin etmek için kaç ölçüm gerekir? Örneklem büyüklüğü iki katına çıkarılırsa hassasiyet nasıl değişir?
\end{alistirma}
\alcozum{%
$n\geq(1.96\cdot5/0.5)^2=384.16\Rightarrow n=385$. $n=770$'de: ME$=1.96\cdot5/\sqrt{770}=0.354$ mm, yani $1/\sqrt{2}\approx0.707$ katına düşer. Hassasiyet $\sqrt{2}$ faktörüyle artmış.%
}

\begin{alistirma}[Al.13.7]{C}{Pivot yöntemi — Uniform}
$X_1,\ldots,X_n\iid\Uniform[0,\theta]$. $Q=X_{(n)}/\theta$'nın dağılımını bul ve $\theta$ için $(1-\alpha)$ güven aralığını türet.
\end{alistirma}
\alcozum{%
$F_{X_{(n)}}(t)=(t/\theta)^n$; dolayısıyla $Q=X_{(n)}/\theta\sim\mathrm{Beta}(n,1)$ (KDF: $F_Q(q)=q^n$). $\Prob(q_{L}\leq Q\leq1)=1-\alpha$ ile $q_L=\alpha^{1/n}$. Çöz: $X_{(n)}/1\leq\theta$ her zaman; alt sınır $\Prob(Q\geq q_L)=1-q_L^n=1-\alpha$. GA: $\theta\in[X_{(n)},\;X_{(n)}/\alpha^{1/n}]$.%
}

\begin{alistirma}[Al.13.8]{B}{Eşzamanlı güven aralığı}
3 parametre $(\theta_1,\theta_2,\theta_3)$ için bireysel \%95 aralıklar kur isteniyor. (a) Bonferroni düzeltmesiyle her aralığın hangi güven düzeyinde kurulması gerektiğini söyle. (b) Eğer her aralığın uzunluğu 4 ise ve Bonferroni uygulanırsa FWER en az $1-3\cdot(1-0.9833)=0.95$ olduğunu göster.
\end{alistirma}
\alcozum{%
(a) Bonferroni: $\alpha_j=0.05/3\approx0.0167$; her aralık $\%98.33$ güven düzeyinde. (b) Her aralık hata $\Prob(A_j)=0.0167$. FWER $\leq\sum_j\Prob(A_j)=3\cdot0.0167=0.05$; yani toplam hata $\leq\alpha=0.05$. Bu doğrudan Birleşim sınırı argümanıdır.%
}

\begin{alistirma}[Al.13.9]{B}{En kısa GA karşılaştırması}
$n=10$, $s^2=9.6$, Normal popülasyon. (a) $\sigma^2$ için eşit kuyruk \%90 GA'sını hesapla ($\chi^2_{9,0.05}=16.92$, $\chi^2_{9,0.95}=3.33$). (b) Üst tek taraflı \%90 GA'sını kur. (c) Hangi GA daha kısa?
\end{alistirma}
\alcozum{%
(a) Eşit kuyruk: $[9\cdot9.6/16.92,\;9\cdot9.6/3.33]=[5.11,\;25.95]$; uzunluk $=20.84$. (b) Tek taraflı üst: $(-\infty,\;9\cdot9.6/\chi^2_{9,0.90})=(-\infty,\;9\cdot9.6/4.17)=(-\infty,\;20.72)$; tek yönlü GA daha kısa ama üst sınırı yoktur. Eşit kuyruk iki taraflı GA burada uzundur çünkü $\chi^2$ asimetriktir; en kısa iki taraflı GA sayısal optimizasyon gerektirir.%
}

\begin{alistirma}[Al.13.10]{C}{Profil olabilirlik GA}
$X_1,\ldots,X_n\iid\Normal(\mu,\sigma^2)$, ikisi de bilinmiyor. $\psi=\mu$ ilgi parametresi. (a) Profil log-olabilirliği $\ell_P(\mu)$'yü hesapla. (b) $R_P(\mu)=2[\ell(\hat\mu,\hat\sigma^2)-\ell_P(\mu)]$ ifadesini hesapla ve $n(\bar X-\mu)^2/S^2$ ile ilişkilendir.
\end{alistirma}
\alcozum{%
(a) $\ell(\mu,\sigma^2)=-\frac{n}{2}\log\sigma^2-\frac{1}{2\sigma^2}\sum(X_i-\mu)^2$. $\hat\sigma^2(\mu)=\frac1n\sum(X_i-\mu)^2$. $\ell_P(\mu)=-\frac{n}{2}\log\hat\sigma^2(\mu)-n/2$. (b) $R_P(\mu)=n\log(\hat\sigma^2(\mu)/\hat\sigma^2)=n\log(1+t^2/(n-1))$ ($t=\sqrt{n}(\bar X-\mu)/S$). $R_P\leq\chi^2_{1,\alpha}$ koşulu $|t|\leq t_{n-1,\alpha/2}$'ye indirgenir — profil olabilirlik GA, $t$-aralığıyla özdeştir.%
}
